\chapter{Quality Assurance: Verification and Validation}
\label{ch:quality}

This chapter documents the quality-assurance programme of the Smart Home Energy
Optimizer (SHEO). It establishes the conceptual vocabulary of software quality,
positions the project's testing effort within the classical V-model, and then
demonstrates \emph{both} families of dynamic testing required of a rigorous
verification campaign: black-box (specification-based) testing and white-box
(structure-based) testing. The discussion is grounded in the project's
\textbf{65 automated tests} and concludes with a requirements-to-tests
traceability summary and a forward-looking classification of anticipated
maintenance.

\section{Software Quality, Verification and Validation}
\label{sec:qa-vv}

Following the IEEE Standard Glossary of Software Engineering Terminology
(IEEE~610.12)~\cite{ieee610}, \emph{software quality} is the degree to which a system meets
specified requirements and, in turn, the degree to which it meets stated and
implied user needs and expectations. Quality is therefore not a single attribute
but a composite of functional correctness together with non-functional
properties such as reliability, security, safety and maintainability. For SHEO
the load-bearing non-functional families are safety (NFR-SAF, protecting
safety-critical appliances and rate-limiting the air-conditioner compressor) and
security (NFR-SEC, authentication, role-based access control and never storing
plaintext secrets).

Two complementary activities establish quality. \emph{Verification} asks
``are we building the product right?'' --- it is the developer-facing check that
each artefact conforms to its specification and to the artefact that preceded it
(requirements, design, code). \emph{Validation} asks ``are we building the right
product?'' --- it is the user-facing check that the delivered system actually
satisfies the stakeholder's real needs. In SHEO, verification is realised by the
unit and integration tests that confirm the implementation honours its
contracts, while validation is realised by the system and acceptance tests that
exercise the application through its public interface against the acceptance
criteria agreed with the customer --- here the SHEO stakeholder in the IEEE
sense, namely the Administrator who commissions and accepts the system. The two
are not interchangeable: a system can
be verified (faithful to its specification) yet invalid (the specification was
wrong), which is precisely why both arms of the V-model below are required.

\section{The V-model}
\label{sec:qa-vmodel}

The \emph{V-model} arranges the software life-cycle so that every constructive
(development) phase on the descending left arm is paired with a corresponding
destructive (test) phase on the ascending right arm. The horizontal links express
the principle that the test cases for a level are derived from, and validate
against, the artefact produced by the phase opposite it: acceptance tests are
written against the requirements, system tests against the architecture, and
integration tests against the detailed design, and so on.
Figure~\ref{fig:vmodel} draws the model and maps SHEO's concrete test
levels onto it.

\begin{figure}[H]
\centering
\resizebox{\linewidth}{!}{%
\begin{tikzpicture}[
    node distance=6mm,
    box/.style={draw, rounded corners, align=center, minimum height=8mm,
                minimum width=30mm, font=\small},
    dev/.style={box, fill=blue!8},
    test/.style={box, fill=green!10},
    link/.style={dashed, -{Latex[length=2mm]}, gray},
    flow/.style={-{Latex[length=2mm]}, thick}
]
% Left (development) arm, descending
\node[dev] (req)  {Requirements\\analysis};
\node[dev] (arch) [below=of req, xshift=10mm]  {Architectural\\design};
\node[dev] (dd)   [below=of arch, xshift=10mm] {Detailed\\design};
\node[dev] (code) [below=of dd, xshift=10mm]   {Coding};
% Right (test) arm, ascending
\node[test] (acc) [right=42mm of req]  {Acceptance\\testing};
\node[test] (sys) [below=of acc, xshift=-10mm]  {System\\testing};
\node[test] (int) [below=of sys, xshift=-10mm]  {Integration\\testing};
\node[test] (unit)[below=of int, xshift=-10mm]  {Unit\\testing};
% Development flow
\draw[flow] (req)  -- (arch);
\draw[flow] (arch) -- (dd);
\draw[flow] (dd)   -- (code);
% Coding to unit testing (bottom of the V)
\draw[flow] (code) -- (unit);
% Test flow (ascending)
\draw[flow] (unit) -- (int);
\draw[flow] (int)  -- (sys);
\draw[flow] (sys)  -- (acc);
% Horizontal verification/validation links
\draw[link] (req)  -- (acc) node[midway, above, font=\scriptsize] {validates};
\draw[link] (arch) -- (sys) node[midway, above, font=\scriptsize] {verifies};
\draw[link] (dd)   -- (int) node[midway, above, font=\scriptsize] {verifies};
\end{tikzpicture}}
\caption{The V-model mapped onto SHEO. Each development phase (left, shaded blue)
is checked by the test level opposite it (right, shaded green): the topmost link
\emph{validates} the requirements against the stakeholder's real needs, while the
lower links \emph{verify} each design artefact against its specification. Dashed
links denote the artefact each test level is derived from.}
\label{fig:vmodel}
\end{figure}

SHEO populates every level of the right arm. \emph{Unit testing} exercises the
pure domain functions in isolation; \emph{integration testing} drives the wired
stack (API router $\rightarrow$ service $\rightarrow$ repository $\rightarrow$
persistence) across real boundaries; \emph{system testing} asserts aggregate,
user-visible outputs end-to-end on a fully seeded application; and
\emph{acceptance testing} demonstrates the customer's acceptance focus items.
The sequence-level interactions that integration tests confirm are illustrated in
the design chapter (Figure~\ref{fig:seq-control}).

\section{Test Strategy and Levels}
\label{sec:qa-strategy}

The automated suite comprises \textbf{65 tests}. It runs in-process against the
application's public interface using an in-memory test harness, requiring no live
server, network or hardware; a disposable database is recreated for each run, so
the suite is fully deterministic. The split between levels is physical: pure-logic
unit tests touch no transport layer and no database, whereas every other suite
boots a fully seeded application and drives it through its public endpoints.

Each test is expressed in the canonical test-case form
$\langle$input, expected outcome$\rangle$: an input (or input sequence) is applied
and the actual result is compared against an \emph{expected output that is defined
in advance}, never inferred after observing the program. A collection of such
cases that share a fixture and a target constitutes a \emph{test suite} --- the
project organises its tests into nine such suites, each devoted to one
cross-cutting concern. The strategy observes three classical testing principles:

\begin{itemize}[leftmargin=1.4em]
  \item the \textbf{tester $\ne$ developer} separation of concern: tests are
        written from the specification (the requirements and the capability
        contract) rather than to ratify whatever the code happens to do;
  \item the \textbf{expected output is fixed in advance}, so a test can only pass
        by the implementation being correct, not by redefining correctness;
  \item \textbf{regression}: because the suite is automated and deterministic, it
        is re-run on every change to guard previously verified behaviour against
        accidental breakage.
\end{itemize}

Table~\ref{tab:levels} summarises how the 65 tests distribute across the four
levels.

\begin{table}[htbp]
\centering
\caption{Distribution of the 65 automated tests across V-model levels.}
\label{tab:levels}
\begin{tabularx}{\textwidth}{@{}l L{0.52\textwidth} Y@{}}
\toprule
\textbf{Level} & \textbf{What it checks} & \textbf{Representative target} \\
\midrule
Unit        & Pure domain functions in isolation: capability validation,
              tariff pricing, billing-cycle and time utilities, simulator energy
              model, password hashing. & the pure-domain logic suite \\
Integration & Multiple layers cooperating across a real boundary: the wired
              request path through service, repository and persistence, plus the
              event bus and adapter wiring. & the device-capability and rules
              suites \\
System      & End-to-end behaviour of a fully seeded application through public
              endpoints: dashboards, reports, recommendations, savings. &
              the monitoring and savings suites \\
Acceptance  & The customer's acceptance focus items, realised by the union of the
              system and integration tests against the demonstration seed. & (see
              Section~\ref{sec:qa-results}) \\
\bottomrule
\end{tabularx}
\end{table}

\section{Black-box Testing}
\label{sec:qa-blackbox}

Black-box (specification-based) testing derives test cases from the external
specification of a component without reference to its internal structure. The
suite deliberately applies three systematic black-box design techniques rather
than choosing inputs ad hoc.

\subsection{Equivalence Partitioning}
\label{sec:qa-ep}

\emph{Equivalence partitioning} divides the input domain into classes within
which the program is expected to behave identically, so that one representative
per class suffices. The air-conditioner target-temperature control accepts an
integer temperature in the closed range $[16, 30]$ degrees Celsius. Its domain
partitions into three classes:

\begin{itemize}[leftmargin=1.4em]
  \item a \textbf{valid} class --- any value inside $[16, 30]$ (representative:
        $23$), expected to be accepted;
  \item an \textbf{invalid-low} class --- values below $16$, expected to be
        rejected;
  \item an \textbf{invalid-high} class --- values above $30$, expected to be
        rejected.
\end{itemize}

The principle generalises beyond numeric ranges. The enumerated-and-toggle
partitioning test divides such controls into valid representatives (a valid mode
and power setting) and invalid ones (an invalid mode; an invalid power value),
while a companion test treats an unsupported control (for instance a brightness
setting directed at a device that has none) as its own invalid class.

\subsection{Boundary Value Analysis}
\label{sec:qa-bva}

\emph{Boundary value analysis} (BVA) exploits the empirical observation that
faults cluster at the edges of equivalence classes. For the $[16, 30]$ range it
prescribes testing each boundary, its immediate neighbour just outside, and a
nominal interior value. Table~\ref{tab:bva} gives the canonical BVA set, plus an
out-of-domain negative value; all six cases are run as a single parametrised
boundary test driven directly against the capability-validation routine.

\begin{table}[htbp]
\centering
\caption{Boundary value analysis for the air-conditioner target range $[16, 30]$.}
\label{tab:bva}
\begin{tabular}{@{}c l c@{}}
\toprule
\textbf{Input} & \textbf{Class} & \textbf{Expected} \\
\midrule
$16$  & lower boundary (valid)            & accept \\
$30$  & upper boundary (valid)            & accept \\
$23$  & nominal interior                  & accept \\
$15$  & just below lower boundary         & reject \\
$31$  & just above upper boundary         & reject \\
$-5$  & far invalid (negative)            & reject \\
\bottomrule
\end{tabular}
\end{table}

This realises the canonical ``min, $\text{min}-1$, max, $\text{max}+1$,
nominal'' boundary set: the value $16$ exercises the $\ge \text{min}$ guard, $15$
its lower neighbour, $30$ the $\le \text{max}$ guard, and $31$ its upper
neighbour.

\subsection{Decision-table Testing}
\label{sec:qa-dt}

When an outcome depends on a \emph{combination} of conditions, a
\emph{decision table} enumerates the relevant condition combinations (rules) and
their required actions, ensuring that no combination is overlooked. The outcome
of a device command depends on whether the device is online, whether the value is
in range, and whether the command is suppressed for safety (a safety-critical
auto power-off, or a rate-limited compressor command). The resulting outcomes are
\textsc{success}, \textsc{rejected}, \textsc{timeout} and \textsc{skipped}.
Table~\ref{tab:decision} captures the rules and describes the test that exercises
each.

\begin{table}[htbp]
\centering
\caption{Decision table for command outcomes. ``--'' denotes a condition whose
value is irrelevant once an earlier condition has determined the outcome.}
\label{tab:decision}
\begin{tabularx}{\textwidth}{@{}c c c >{\centering\arraybackslash}p{0.16\textwidth} L{0.34\textwidth}@{}}
\toprule
\textbf{Online?} & \textbf{In range?} & \textbf{Safety-blocked?} & \textbf{Outcome} & \textbf{Covering test} \\
\midrule
no  & --  & --  & \textsc{timeout}  & the offline-command timeout test \\
yes & no  & --  & \textsc{rejected} & the out-of-range rejection test \\
yes & yes & yes & \textsc{skipped}  & the safety-critical auto-power-off suppression test \\
yes & yes & yes & \textsc{rejected} & the compressor rate-limit test \\
yes & yes & no  & \textsc{success}  & the in-range acceptance test \\
\bottomrule
\end{tabularx}
\end{table}

A related decision table over the conditions (temperature change?),
(safety-critical tag?) and (unattended power-off?) governs the validation
\emph{warnings} surfaced to the user, and is exercised by the safety-warning
validation test.

\section{White-box Testing}
\label{sec:qa-whitebox}

White-box (structure-based) testing derives test cases from the internal control
structure of the code, with the goal of executing its statements, branches and
paths. We analyse the range-validation logic of the capability-validation routine
--- the procedure that enforces the $[16, 30]$ contract tested above. For a
range-typed control it proceeds through a fixed sequence of guarded steps: it
first coerces the supplied value to a number and rejects it if that is not
possible; it then rejects any value below the declared minimum, and any value
above the declared maximum; for a control declared with an integer step it
rounds the value to a whole number; and finally it accepts the normalised value.
The corresponding \emph{control flow graph} (CFG), in which decisions are drawn
as diamonds and basic blocks as rectangles, is given in Figure~\ref{fig:cfg};
its numbered nodes correspond to these steps.

\begin{figure}[H]
\centering
\begin{tikzpicture}[
    node distance=8mm and 16mm,
    block/.style={draw, rounded corners, align=center, font=\scriptsize,
                  minimum width=20mm, minimum height=7mm},
    dec/.style={draw, diamond, aspect=2, align=center, font=\scriptsize,
                inner sep=1pt},
    term/.style={draw, ellipse, align=center, font=\scriptsize, fill=red!8},
    ok/.style={draw, ellipse, align=center, font=\scriptsize, fill=green!12},
    edge/.style={-{Latex[length=2mm]}}
]
\node[block] (n1) {1: float(value)};
\node[dec]   (n1d) [below=of n1] {numeric?};
\node[term]  (n2)  [right=of n1d] {2: reject\\``numeric''};
\node[dec]   (n3)  [below=of n1d] {3: num\\$<$ min?};
\node[term]  (n4)  [right=of n3] {4: reject\\``$\ge$ min''};
\node[dec]   (n5)  [below=of n3] {5: num\\$>$ max?};
\node[term]  (n6)  [right=of n5] {6: reject\\``$\le$ max''};
\node[dec]   (n7)  [below=of n5] {7: step?};
\node[block] (n8)  [right=of n7] {8: round};
\node[ok]    (n9)  [below=of n7] {9: accept};

\draw[edge] (n1)  -- (n1d);
\draw[edge] (n1d) -- node[above,font=\tiny]{no} (n2);
\draw[edge] (n1d) -- node[right,font=\tiny]{yes} (n3);
\draw[edge] (n3)  -- node[above,font=\tiny]{yes} (n4);
\draw[edge] (n3)  -- node[right,font=\tiny]{no} (n5);
\draw[edge] (n5)  -- node[above,font=\tiny]{yes} (n6);
\draw[edge] (n5)  -- node[right,font=\tiny]{no} (n7);
\draw[edge] (n7)  -- node[above,font=\tiny]{yes} (n8);
\draw[edge] (n7)  -- node[right,font=\tiny]{no} (n9);
\draw[edge] (n8)  |- (n9);
\end{tikzpicture}
\caption{Control flow graph of the range-validation branch of the
capability-validation routine. Diamonds are decisions; red terminals are
rejections; the green terminal is the accepting return. Numbers correspond to the
guarded steps described above.}
\label{fig:cfg}
\end{figure}

The CFG exposes four decisions (nodes 1, 3, 5, 7) and therefore eight branch
outcomes. Table~\ref{tab:coverage} gives a small set of
$\langle$input $\rightarrow$ expected$\rangle$ cases that, taken together, achieve
\emph{path} coverage of this subgraph (every feasible execution path), which
subsumes both branch and statement coverage.

\begin{table}[htbp]
\centering
\caption{White-box coverage cases for the range branch (air-conditioner target,
with an integer step of $1$). The node sequence shows the path taken.}
\label{tab:coverage}
\begin{tabularx}{\textwidth}{@{}l c L{0.30\textwidth} Y@{}}
\toprule
\textbf{Input} & \textbf{Expected} & \textbf{Path (nodes)} & \textbf{Exercises} \\
\midrule
non-numeric    & reject (not numeric) & $1 \rightarrow 2$           & node-1 false branch \\
$15$           & reject ($<$ min)     & $1 \rightarrow 3 \rightarrow 4$ & node-3 true branch \\
$31$           & reject ($>$ max)     & $1 \rightarrow 3 \rightarrow 5 \rightarrow 6$ & node-5 true branch \\
$23$           & accept (rounded)     & $1 \rightarrow 3 \rightarrow 5 \rightarrow 7 \rightarrow 8 \rightarrow 9$ & node-7 true, normalising path \\
\bottomrule
\end{tabularx}
\end{table}

The four cases above execute all eight branch outcomes (node 7's false branch is
taken by any range control declared without an integer step, e.g. a continuous
control), and in doing so visit every statement node, so they satisfy branch and
statement coverage as well. This illustrates the standard \emph{subsumption}
ordering of coverage strength:
\[
  \text{path coverage} \;\ge\; \text{branch coverage} \;\ge\; \text{statement coverage},
\]
that is, achieving a stronger criterion guarantees the weaker ones but not
conversely. The boundary inputs $15$, $23$ and $31$ from Table~\ref{tab:bva} thus
do double duty: chosen by black-box BVA, they also drive the most fault-prone
branches of the white-box CFG.

\section{Results and Traceability}
\label{sec:qa-results}

The full suite passes: \textbf{65 tests, 65 passed}. Because tests are derived
from the specification, every requirement family is traceable to at least one
executable check, and conversely every test traces back to a requirement.
Table~\ref{tab:trace} summarises the mapping from requirement families to
representative covering tests.

\begin{table}[htbp]
\centering
\caption{Requirements-to-tests traceability (representative tests).}
\label{tab:trace}
\begin{tabularx}{\textwidth}{@{}l L{0.32\textwidth} Y@{}}
\toprule
\textbf{Requirement family} & \textbf{Concern} & \textbf{Behaviour verified} \\
\midrule
REQ-4.1 & Monitoring (live power, cumulative energy, offline detection) &
  that the home dashboard reports the expected metrics, and that a command to an
  offline device times out \\
REQ-4.2 & Device control via capability schema (load, validate, command) &
  that the capability schema drives the available controls, and that an
  out-of-range command is rejected \\
REQ-4.3 & Rules, conflict detection, auto-action opt-in and undo &
  that conflicting rules are detected, and that automatic actions are disabled
  by default \\
REQ-4.4 & Recommendation generation, ranking and approval &
  that an accepted recommendation becomes a rule, and that recommendations are
  ranked and capped \\
REQ-4.5 & Bill-saving estimation and cycle reporting &
  that the saving estimate follows the specified formula, and that the savings
  summary reports the billing cycle \\
NFR-SAF & Compressor rate limit, safety-critical protection, warnings &
  the compressor rate-limit test, and the safety-critical auto-power-off
  suppression test \\
NFR-SEC & Authentication, RBAC, no plaintext secrets &
  that an unauthenticated request is rejected, and that a password survives a
  hash round-trip without ever being stored in plaintext \\
\bottomrule
\end{tabularx}
\end{table}

Both styles of testing are present and traceable: pure-logic unit tests pin the
algorithms (validation, pricing, the simulator energy differential), while seeded
in-process tests driven through the public interface prove the wired system
behaves correctly. Negative and boundary paths are first-class throughout: rejected,
timeout and skipped command outcomes, out-of-range and unsupported controls, bad
passwords, unauthenticated access and forbidden cross-role actions are all
explicitly exercised. The principal residual gaps --- direct assertion of
notification delivery, frontend rendering, and non-functional load and transport
security --- are documented as known limitations to be closed with a notification
endpoint test, manual UI acceptance, and a dedicated load and TLS harness
respectively.

\section{Maintenance}
\label{sec:qa-maintenance}

Verification and validation do not end at delivery; they continue through the
Support phase of the life-cycle. Anticipated changes to SHEO classify into the
four standard categories of software maintenance, each of which would re-enter the
regression suite before release:

\begin{itemize}[leftmargin=1.4em]
  \item \textbf{Corrective} --- repairing a latent fault. For example, fixing an
        off-by-one error should the air-conditioner target guard ever admit $31$;
        the corresponding boundary test in Table~\ref{tab:bva} would reproduce
        and then guard the fix.
  \item \textbf{Adaptive} --- accommodating a changed environment. For example,
        updating the tariff math when the electricity utility revises its tiered
        price schedule, re-validated by the tariff-pricing unit tests.
  \item \textbf{Perfective} --- improving function or quality without changing
        external behaviour. For example, adding a new device type (a water
        heater) by declaring one capability schema entry and one adapter, then
        extending the capability and command tests to cover it.
  \item \textbf{Preventive} --- restructuring to forestall future defects. For
        example, hardening the rate-limiter against clock skew before any failure
        is observed, protected by the existing rate-limit regression test.
\end{itemize}

This classification realises the Support phase of the life-cycle introduced in
the process chapter, and closes the verification-and-validation loop: each kind
of change is admitted only once the relevant tests --- old and new --- pass,
keeping the requirements-to-tests traceability of Table~\ref{tab:trace} honest as
the system evolves.
